% presentation_organisme_encadrant.tex
\documentclass[12pt,a4paper]{article}
\usepackage[utf8]{inputenc}
\usepackage[T1]{fontenc}
\usepackage[french]{babel}
\usepackage{geometry}
\usepackage{setspace}
\usepackage{parskip} % pour espacement entre paragraphes

\geometry{top=2.5cm, bottom=2.5cm, left=2.5cm, right=2.5cm}
\setstretch{1.3}

\begin{document}

\section*{Présentation de l’organisme d’accueil et de l’encadrant}

Le projet a été réalisé au sein de la \textbf{Faculté des Sciences de Tunis (FST)}, située à Tunis, en Tunisie. La FST est un établissement \textbf{public} rattaché à l’Université de Tunis El Manar, spécialisé dans l’enseignement supérieur et la recherche scientifique dans les domaines des sciences fondamentales et appliquées.

La faculté est structurée en plusieurs départements, parmi lesquels le \textbf{Département d’Informatique}, qui regroupe des enseignants-chercheurs, des laboratoires de recherche et des plateformes techniques dédiées aux projets académiques. Le projet s’est déroulé dans ce cadre, en lien direct avec les enseignements du module \textbf{Big Data}, permettant d’appliquer les concepts théoriques à un projet concret dans le domaine de la finance et de la détection de fraude bancaire.

L’encadrement du projet a été assuré par \textbf{Mme Manel Zekri}, enseignante-chercheure au Département d’Informatique de la FST. Elle a supervisé l’ensemble des étapes du projet, depuis la définition du sujet jusqu’à la validation des résultats, en apportant un suivi méthodologique et technique. Son rôle a été essentiel pour orienter le travail, corriger les difficultés rencontrées et assurer la qualité scientifique et pratique du projet.

\end{document}
